\chapter{Monad and Algebraic Laws}
\label{app:laws}

\chapterorient{The combinators of \cref{ch:combinators} and the state monad of
\cref{ch:monad} are only as trustworthy as the rewrites they license. This appendix
is the reference catalogue of those rewrites: the monad laws that let you refactor a
\code{do}-block without changing what it computes, the state-effect algebra of
\code{get}/\code{put}/\code{modify}, and the equational laws of the tensor
primitives. But the catalogue has a second half that is every bit as important as the
first, and far easier to forget---the laws that do \emph{not} hold. A refactor that
\emph{looks} valid because it resembles a law you half-remember can silently change
the algorithm, or, more insidiously, change only its last bits. Those silent traps
are pinned here, each with the reason it fails and the harm it does. Read the
right-hand column as carefully as the left.}

\section{How to read this appendix}
\label{sec:cread}

The chapters of Parts~III and~IV state laws in passing, as they need them. This
appendix gathers them into one place and adds the precision a reference demands: each
law is an \emph{equation} (in math mode) or a \emph{rewrite} (a pseudocode
\code{==}), together with one sentence on what it \emph{guarantees}---what you are
allowed to do because it holds. The companion appendix on the formal calculus
(\cref{app:calculus}) gives the small-step reduction rules these laws sit on top of;
where a law is just a reduction rule renamed, we say so and point there.

Three kinds of statement appear, and the typography distinguishes them.

\begin{description}[style=nextline, leftmargin=0pt, font=\normalfont\itshape]
  \item[A law.] An equation that holds and may be used as a rewrite in either
  direction. Stated in a \code{theorem}, \code{lemma}, or a numbered list.
  \item[A non-law.] An equation that \emph{looks} like a law but is \emph{false}, set
  off explicitly. Each carries the reason it fails and the concrete damage assuming
  it would do. These are collected in \cref{sec:nonlaws} and are the most valuable
  content here.
  \item[A floating-point caveat.] A law that holds over an idealized exact field but
  only \emph{approximately} under IEEE-754 arithmetic. These are a special, sharp
  class of non-law---the equation is mathematically true and numerically false---and
  they are why \cref{app:numerics} exists. They are flagged inline and revisited in
  \cref{sec:cfp}.
\end{description}

\begin{keyidea}
The non-laws are as important as the laws. A law tells you a rewrite is \emph{safe};
a non-law tells you a rewrite that \emph{looks} safe is not. The dangerous case is
not the rewrite that obviously breaks---the types catch that---but the one that
resembles a law you half-remember: hoisting a loop predicate, merging two folds,
reordering a reduction. Such a rewrite can leave the program type-correct, still
running, and computing the wrong thing---or, worse, computing a \emph{bit-different}
thing that diverges only after thousands of iterations. \textbf{This is why a
refactor that looks valid can silently change the last bits.} The right-hand column
of this appendix---the laws that do \emph{not} hold---is the one you consult before
touching working solver code.
\end{keyidea}

\section{The monad laws}
\label{sec:cmonad}

We begin with the three laws that govern \code{>>=} (``bind'') and \code{return}, the
two operations underneath every \code{do}-block of \cref{ch:monad}. A \code{do}-block
is sugar: each line is desugared into a \code{>>=} that threads the state from one
line into the next (\cref{app:calculus} gives the desugaring in full). The three
monad laws are what make that sugar \emph{sound}---they are the guarantee that
rearranging the block in the ways the laws permit changes nothing about what it
computes. They hold for the \code{Solve} monad of \cref{ch:monad} exactly as they
hold for every monad.

\begin{theorem}[The three monad laws]
\label{thm:monad}
For all values \code{a}, recipes \code{m}, and recipe-valued functions \code{f},
\code{g}:
\begin{align}
  \textsf{return}\ a \,\texttt{>>=}\, f
    &\;=\; f\ a
    & \text{(left identity)} \label{eq:leftid}\\[2pt]
  m \,\texttt{>>=}\, \textsf{return}
    &\;=\; m
    & \text{(right identity)} \label{eq:rightid}\\[2pt]
  (m \,\texttt{>>=}\, f) \,\texttt{>>=}\, g
    &\;=\; m \,\texttt{>>=}\, (\lambda x.\, f\ x \,\texttt{>>=}\, g)
    & \text{(associativity)} \label{eq:assoc}
\end{align}
\end{theorem}

\noindent Take them one at a time, each with the refactor it licenses.

\paragraph{Left identity \eqref{eq:leftid}: \code{return} then bind is just
substitution.} If you lift a pure value into a recipe with \code{return} and
immediately feed it to \code{f}, you may as well have called \code{f} on the value
directly. In \code{do}-notation: a line
\lstinline[style=l4style]|x <- return a| followed by a use of \code{x} is identical
to binding \code{x} with a plain \lstinline[style=l4style]|let x = a|. The guarantee:
\emph{a \code{return} immediately consumed adds nothing}---you may always erase it and
substitute. This is the law that lets the over-monadifying mistake of
\cref{ch:monad}'s pitfall be \emph{undone}: a value needlessly wrapped in
\code{return} and then bound back out collapses to the pure binding it should have
been.

\paragraph{Right identity \eqref{eq:rightid}: binding into \code{return} is a
no-op.} If a recipe \code{m} runs and you immediately \code{return} whatever it
yielded, changing nothing, the trailing \code{return} can be dropped: the result is
just \code{m}. The guarantee: \emph{a recipe that ends by returning its own result is
that recipe}. This is what lets a \code{do}-block's final
\lstinline[style=l4style]|x <- step; return x| be shortened to \code{step}---the
redundant last line is provably inert.

\paragraph{Associativity \eqref{eq:assoc}: grouping of binds does not matter.} This
is the load-bearing one. It says that when three recipes run in sequence, it makes no
difference whether you parenthesize the sequencing as ``(first two) then third'' or
``first then (last two).'' In \code{do}-notation this is exactly the statement that
\emph{flattening and splitting \code{do}-blocks is safe}: a block

\begin{lstlisting}[style=l4style]
do { a <- m; do { b <- f a; g b } }   -- inner do-block nested
\end{lstlisting}

\noindent computes the same final state and value as the flattened

\begin{lstlisting}[style=l4style]
do { a <- m; b <- f a; g b }          -- flattened
\end{lstlisting}

\noindent because the two are \eqref{eq:assoc} read left-to-right and right-to-left.
The guarantee: \emph{you may regroup, factor out, and inline sub-blocks freely}. When
\cref{ch:monad} factors a Krylov step's body into a helper recipe and splices it back
into the loop, associativity is the warrant that the splice changed nothing.
\Cref{fig:monadlaw} draws this rewrite as a before/after.

\begin{figure}[t]
\centering
\begin{tikzpicture}[font=\footnotesize]
  % ===== LEFT: nested do-block =====
  \begin{scope}
    \node[font=\small\bfseries, text=simBlue] at (0,2.55) {before: nested};
    \node[draw=simGray!55, fill=simBgGray, rounded corners=4pt, thick,
          minimum width=42mm, minimum height=34mm] (outer) at (0,0.4) {};
    \node[anchor=north west, font=\ttfamily\scriptsize] at ($(outer.north west)+(2mm,-2mm)$)
      {do \{};
    \node[anchor=west, font=\ttfamily\scriptsize] at ($(outer.north west)+(4mm,-7mm)$)
      {a <- m;};
    \node[draw=simBlue!60, fill=simBgBlue, rounded corners=3pt, thick,
          minimum width=30mm, minimum height=15mm] (inner) at (0.6,-0.55) {};
    \node[anchor=north west, font=\ttfamily\scriptsize] at ($(inner.north west)+(1.5mm,-1.5mm)$)
      {do \{};
    \node[anchor=west, font=\ttfamily\scriptsize] at ($(inner.north west)+(3.5mm,-6mm)$)
      {b <- f a;};
    \node[anchor=west, font=\ttfamily\scriptsize] at ($(inner.north west)+(3.5mm,-10mm)$)
      {g b \} \}};
  \end{scope}
  % ===== arrow =====
  \node[font=\large, text=simRust] (arr) at (3.3,0.4) {$\equiv$};
  \node[font=\scriptsize\itshape, text=simRust, align=center] at (3.3,-0.55)
    {assoc.\\\eqref{eq:assoc}};
  % ===== RIGHT: flattened =====
  \begin{scope}[xshift=6.6cm]
    \node[font=\small\bfseries, text=simGreen] at (0,2.55) {after: flattened};
    \node[draw=simGreen!60, fill=simBgGreen, rounded corners=4pt, thick,
          minimum width=42mm, minimum height=34mm] (flat) at (0,0.4) {};
    \node[anchor=north west, font=\ttfamily\scriptsize] at ($(flat.north west)+(2mm,-2mm)$)
      {do \{};
    \node[anchor=west, font=\ttfamily\scriptsize] at ($(flat.north west)+(4mm,-8mm)$)
      {a <- m;};
    \node[anchor=west, font=\ttfamily\scriptsize] at ($(flat.north west)+(4mm,-14mm)$)
      {b <- f a;};
    \node[anchor=west, font=\ttfamily\scriptsize] at ($(flat.north west)+(4mm,-20mm)$)
      {g b \}};
  \end{scope}
  \node[font=\scriptsize\itshape, text=simGray, align=center] at (3.3,-2.4)
    {same threaded state, same yielded value};
\end{tikzpicture}
\caption[A monad law as a \texttt{do}-block rewrite]{Associativity \eqref{eq:assoc}
read as a \code{do}-block refactor. The nested block on the left---an inner
\code{do} spliced into an outer one---computes exactly the same final state and
yielded value as the flattened block on the right. The two are the two
parenthesizations of \eqref{eq:assoc}: \lstinline[style=l4style]|(m >>= f) >>= g|
versus \lstinline[style=l4style]|m >>= (\x -> f x >>= g)|. Because the law holds, you
may flatten, factor, and inline sub-blocks at will; the threading the
\code{do}-notation hides is rearranged but its net effect is invariant. This is the
formal license behind every ``pull this step out into a helper'' refactor in
\cref{ch:monad}.}
\label{fig:monadlaw}
\end{figure}

\begin{notationbox}
The three laws are stated as \emph{equalities} here, but in \cref{app:calculus} they
appear as left-to-right \emph{reduction rules} (the arrow $\to$). The direction
matters for an evaluator---reduction always moves left to right---but for
\emph{refactoring} both directions are sound, because the relation is an equality.
When this appendix writes a law with \code{=} or $=$ you may rewrite either way; when
\cref{app:calculus} writes the same fact with $\to$ it is fixing the evaluation
direction. Same fact, two uses.
\end{notationbox}

\section{The state-effect laws}
\label{sec:cstate}

The monad laws of \cref{sec:cmonad} hold for \emph{every} monad; they know nothing of
state. The \code{Solve} monad adds three state primitives---\code{get}, \code{put},
\code{modify}---and these obey their own small algebra, the rules for how reads and
writes of the threaded \code{SimState} interact. This algebra is what lets you
simplify a chain of state operations into fewer, clearer ones, and it is the formal
content behind \cref{ch:monad}'s claim that \code{modify} is the one honest verb for
a state change.

\begin{lemma}[The \texttt{get}/\texttt{put}/\texttt{modify} algebra]
\label{lem:state}
Writing \lstinline[style=l4style]|do {...}| for sequencing and $f, g$ for pure
\code{SimState -> SimState} functions:
\begin{align}
  \textsf{modify}\ f
    &\;=\; \textsf{do}\ \{\, s \leftarrow \textsf{get};\ \textsf{put}\ (f\ s) \,\}
    & \text{(\code{modify} \emph{is} get-apply-put)} \label{eq:modifydef}\\[2pt]
  \textsf{modify}\ (f \circ g)
    &\;=\; \textsf{do}\ \{\, \textsf{modify}\ g;\ \textsf{modify}\ f \,\}
    & \text{(\code{modify} fusion)} \label{eq:modifyfuse}\\[2pt]
  \textsf{do}\ \{\, s \leftarrow \textsf{get};\ \textsf{put}\ s \,\}
    &\;=\; \textsf{return}\ ()
    & \text{(get-then-put is inert)} \label{eq:getput}\\[2pt]
  \textsf{do}\ \{\, \textsf{put}\ s_1;\ \textsf{put}\ s_2 \,\}
    &\;=\; \textsf{put}\ s_2
    & \text{(a later \code{put} wins)} \label{eq:putput}\\[2pt]
  \textsf{do}\ \{\, \textsf{put}\ s;\ \textsf{get} \,\}
    &\;=\; \textsf{do}\ \{\, \textsf{put}\ s;\ \textsf{return}\ s \,\}
    & \text{(get after put sees the put)} \label{eq:putget}
\end{align}
\end{lemma}

\noindent These five are the complete working algebra; everything else follows by
combining them with the monad laws. Read each as a rewrite you are permitted.

\paragraph{\code{modify} is get-apply-put \eqref{eq:modifydef}.} This is the
\emph{definition} of \code{modify}, promoted to a law because it justifies both
expansions: any \code{modify} may be unfolded into an explicit get/apply/put, and---
read backwards---any get/apply/put triple may be \emph{folded} into a single
\code{modify}. The backward direction is the useful one: it is how three lines of
state plumbing collapse into the one named verb that, per \cref{ch:monad}, makes the
change stand out. \Cref{ch:monad}'s preference for \code{modify} over \code{put} is
this law's left-hand side being more honest than its right.

\paragraph{\code{modify} fusion \eqref{eq:modifyfuse}.} Two consecutive
\code{modify}s are one \code{modify} by the composed function. The guarantee:
\emph{adjacent state updates merge}. If a step bumps the counter and then, later in
the same block, records a residual---two \code{modify}s---they may be fused into a
single \code{modify} that does both, or, read the other way, one fat \code{modify}
may be split into stages for clarity. Note the order: \eqref{eq:modifyfuse} fuses to
$f \circ g$ with \code{g} applied \emph{first} (it ran first), so the composition
reads right-to-left exactly as function composition does. Getting this order backward
is a classic slip---which is why we write it out.

\paragraph{Get-then-put is inert \eqref{eq:getput}.} Reading the state and writing
back exactly what you read changes nothing: it equals \code{return ()}. The
guarantee: \emph{a read followed by an identity write can be deleted}. This is the
state-effect analog of the monad's right identity \eqref{eq:rightid}, and it is the
law that catches a particular dead-code shape---a refactor that left behind a
\lstinline[style=l4style]|s <- get; put s| with no change in between is doing nothing
and may be removed.

\paragraph{A later \code{put} wins \eqref{eq:putput}.} Two \code{put}s in a row: the
first is overwritten wholesale, so only the second survives. The guarantee:
\emph{a write with no intervening read is dead}. If two \code{put}s sit adjacent with
nothing reading the state between them, the first is pure overhead. (Contrast
\code{modify}: two \code{modify}s do \emph{not} collapse to the last one---each reads
the state the previous wrote, so they compose \eqref{eq:modifyfuse} rather than
overwrite. The difference is exactly that \code{modify} reads and \code{put} does
not, which is the algebra's whole texture.)

\paragraph{Get after put sees the put \eqref{eq:putget}.} If you \code{put} a state
and then \code{get}, the \code{get} yields exactly what you just \code{put}---reads
observe the most recent write. The guarantee: \emph{state is consistent within a
block}; there is no stale-read hazard once the monad threads for you. This is the law
that formalizes the line-by-line trace of \cref{ch:monad}'s walked \code{do}-block,
where the second \code{get} saw the \code{modify}'d state, not the original.

\Cref{fig:state} draws the four interaction laws \eqref{eq:modifyfuse}--\eqref{eq:putget}
as a small interaction table: what each adjacent pair of primitives simplifies to.

\begin{figure}[t]
\centering
\begin{tikzpicture}[font=\footnotesize]
  % header row
  \node[font=\scriptsize\bfseries, text=simGray] at (0,0) {\emph{then} $\downarrow$ / \emph{first} $\rightarrow$};
  \node[opbox, minimum width=16mm] (hg) at (2.2,0) {\code{get}};
  \node[opbox, minimum width=16mm] (hp) at (4.4,0) {\code{put s'}};
  \node[opbox, minimum width=16mm] (hm) at (6.6,0) {\code{modify g}};
  % row 1: get
  \node[opbox, minimum width=16mm] (rg) at (0,-1.0) {\code{get}};
  \node[font=\scriptsize, align=center] at (2.2,-1.0) {(idempotent\\read)};
  \node[font=\scriptsize, text=simRust, align=center] at (4.4,-1.0) {yields \code{s'}\\\eqref{eq:putget}};
  \node[font=\scriptsize, align=center] at (6.6,-1.0) {yields \code{g s}};
  % row 2: put s
  \node[opbox, minimum width=16mm] (rp) at (0,-2.2) {\code{put s}};
  \node[font=\scriptsize, align=center] at (2.2,-2.2) {read then\\overwrite};
  \node[font=\scriptsize, text=simRust, align=center] at (4.4,-2.2) {\code{put s}\\\eqref{eq:putput}};
  \node[font=\scriptsize, align=center] at (6.6,-2.2) {\code{put (g s)}};
  % row 3: modify f
  \node[opbox, minimum width=16mm] (rm) at (0,-3.4) {\code{modify f}};
  \node[font=\scriptsize, align=center] at (2.2,-3.4) {read then\\\code{f}};
  \node[font=\scriptsize, align=center] at (4.4,-3.4) {ignores \code{s'}};
  \node[font=\scriptsize, text=simRust, align=center] at (6.6,-3.4) {\code{modify}\\$(f\!\circ\! g)$ \eqref{eq:modifyfuse}};
  % separators
  \draw[simGray!40] (-1.4,-0.5) -- (7.6,-0.5);
  \draw[simGray!40] (1.1,0.4) -- (1.1,-4.0);
  \node[font=\scriptsize\itshape, text=simGray, align=center] at (3,-4.4)
    {rust cells are the simplification laws of \cref{lem:state};
     \emph{first} primitive runs, \emph{then} the second};
\end{tikzpicture}
\caption[The state-effect algebra]{The \code{get}/\code{put}/\code{modify}
interaction table. Each cell shows what the pair ``run the column primitive, then the
row primitive'' simplifies to. The rust cells are the named laws of
\cref{lem:state}: a \code{get} after a \code{put} yields the put value
\eqref{eq:putget}; a \code{put} after a \code{put} discards the first \eqref{eq:putput};
a \code{modify} after a \code{modify} fuses to the composed function
\eqref{eq:modifyfuse}. The asymmetry between \code{put} (overwrites, so an earlier
one dies) and \code{modify} (reads first, so they compose) is the whole point of the
algebra---and the reason \cref{ch:monad} prefers \code{modify}, the verb that names
its change.}
\label{fig:state}
\end{figure}

\begin{pitfall}
The \code{put}/\code{put} law \eqref{eq:putput} and the \code{modify}/\code{modify}
law \eqref{eq:modifyfuse} look symmetric and are not. Two \code{put}s collapse to the
\emph{last} one---the first is dead. Two \code{modify}s collapse to a
\emph{composition}---neither is dead, because the second \code{modify} reads the state
the first produced. Mistaking one for the other is a real bug: ``simplifying'' two
adjacent \code{modify}s by keeping only the second \emph{drops the first's effect
entirely}. Before collapsing adjacent state writes, ask whether the second
\emph{reads} the state---if it is a \code{modify}, it does, and you must compose, not
discard.
\end{pitfall}

\section{The algebraic operator laws}
\label{sec:coplaws}

We turn from the monad to the tensor primitives it threads. These are the equational
laws of the data-algebra verbs of \cref{ch:combinators}---the BLAS-style operations a
solver step calls---and of the matrix-vector apply. They are the rewrites the
synthesis surface appeals to when arguing that one form of a step computes the same
tensor as another. Unlike the monad laws, several of these carry a floating-point
caveat; those are flagged here and indicted fully in \cref{sec:nonlaws}.

\subsection{The scaling and saxpy laws}

\begin{lemma}[\texttt{axpy} corner laws]
\label{lem:axpy}
For \lstinline[style=l4style]|axpy a x y = a*x + y| (the two-term
\code{linear\_combination} of \cref{ch:combinators}):
\begin{align}
  \textsf{axpy}\ \alpha\ \vv\ \mathbf{0} &\;=\; \alpha \cdot \vv
    & \text{(zero accumulator $\Rightarrow$ scale)} \label{eq:axpyzeroacc}\\
  \textsf{axpy}\ 0\ \vv\ \vy &\;=\; \vy
    & \text{(zero coefficient $\Rightarrow$ identity)} \label{eq:axpyzerocoeff}
\end{align}
\end{lemma}

\noindent These are the two boundary cases of the saxpy. With a zero accumulator the
add drops away and \code{axpy} is a pure scale \eqref{eq:axpyzeroacc}; with a zero
coefficient the scaled term drops away and \code{axpy} returns its accumulator
untouched \eqref{eq:axpyzerocoeff}. The second is the algebraic content of the
implementation's ``skip the term if its coefficient is zero'' branch, lifted to
\code{linear\_combination}'s zero-coefficient term-drop. They license the obvious
peephole simplifications and nothing more.

\subsection{Inner-product laws}

\begin{lemma}[\texttt{dot} and \texttt{inner\_product} laws]
\label{lem:dot}
For the Euclidean inner product \lstinline[style=l4style]|dot v w| and its
operator-weighted generalization
$\inner{\vx}{\vy}_M = \vx^{\!\top} M \vy$:
\begin{align}
  \textsf{dot}\ \vv\ \vw &\;=\; \textsf{dot}\ \vw\ \vv
    & \text{(symmetry, real case)} \label{eq:dotsym}\\
  \textsf{dot}\ (\textsf{axpy}\ \alpha\ \vv\ \vw)\ \vx
    &\;=\; \alpha \cdot \textsf{dot}\ \vv\ \vx + \textsf{dot}\ \vw\ \vx
    & \text{(bilinearity)} \label{eq:dotbilin}\\
  \inner{\vx}{\vx}_M &\;\ge\; 0,\quad
    \inner{\vx}{\vx}_M = 0 \iff \vx = \mathbf{0}
    & \text{(positive-definiteness, SPD $M$)} \label{eq:dotpd}\\
  \inner{\vx}{\vy} &\;=\; \overline{\inner{\vy}{\vx}}
    & \text{(Hermitian symmetry, complex)} \label{eq:dotherm}
\end{align}
\end{lemma}

\noindent Symmetry \eqref{eq:dotsym} lets a solver write \code{dot p Ap} or
\code{dot Ap p} interchangeably. Bilinearity \eqref{eq:dotbilin} is the law that lets
a step-length computation distribute the inner product over a scaled sum---the
workhorse of every Krylov derivation. Positive-definiteness \eqref{eq:dotpd} on an
SPD weight $M$ is the law the norm \code{nrm2} rests its square root on, and the law
that makes the conjugate-gradient energy norm well-defined (\cref{ch:cg}). In the
complex case symmetry weakens to Hermitian symmetry \eqref{eq:dotherm}---the inner
product equals the conjugate of its swap---and the unconjugated bilinear variant
\code{tdot} loses positive-definiteness entirely (it is \emph{not} PSD; see
\cref{sec:nonlaws}).

\begin{pitfall}
Symmetry \eqref{eq:dotsym} is exact in \emph{exact} arithmetic. In floating point,
\code{dot v w} and \code{dot w v} may differ in the last bit if the two
accumulations visit operands in a different order---the multiply is commutative but
the running sum's rounding is order-sensitive. For the symmetric Euclidean dot the
two orders usually coincide, but do not \emph{rely} on bit-identity across a swap;
the safe statement is mathematical equality, not bit-equality. This is the
floating-point caveat of \cref{sec:cfp} in miniature.
\end{pitfall}

\subsection{Matrix-vector apply laws}

\begin{lemma}[\texttt{matvec} laws]
\label{lem:matvec}
For \lstinline[style=l4style]|matvec A v|, the application of a linear operator $A$:
\begin{align}
  \textsf{matvec}\ A\ (\alpha\,\vv)
    &\;=\; \alpha \cdot \textsf{matvec}\ A\ \vv
    & \text{(homogeneity in the vector)} \label{eq:mvhom}\\
  \textsf{matvec}\ (A + B)\ \vv
    &\;=\; \textsf{matvec}\ A\ \vv + \textsf{matvec}\ B\ \vv
    & \text{(additivity in the operator)} \label{eq:mvadd}\\
  (\textsf{matvec}\ A) \circ (\textsf{matvec}\ B)
    &\;=\; \textsf{matvec}\ (A \circ B)
    & \text{(composition)} \label{eq:mvcomp}\\
  \textsf{matvec}\ A\ (\alpha\vv + \beta\vw)
    &\;=\; \alpha\,\textsf{matvec}\ A\ \vv + \beta\,\textsf{matvec}\ A\ \vw
    & \text{(linearity)} \label{eq:mvlin}
\end{align}
\end{lemma}

\noindent Homogeneity \eqref{eq:mvhom} and additivity-in-the-operator \eqref{eq:mvadd}
combine into full linearity \eqref{eq:mvlin}: a \code{matvec} distributes over a
linear combination of its vector argument and pulls scalars out. The
composition law \eqref{eq:mvcomp} is the one with teeth in this book---it says that
applying $B$ then $A$ equals applying the single fused operator $A \circ B$, which is
exactly the algebraic license for the assembled-operator-product forms of the driven
analysis (\cref{ch:combinators} writes the frequency operator
$\mathsf{A}(\omega) = \mathsf{K} + i\omega\,\mathsf{C} - \omega^2\mathsf{M}$ as one
\code{linear\_combination} of operators, justified by \eqref{eq:mvadd}). Additivity
in the operator is also what lets the driven operator be \emph{assembled} as a sum
rather than applied term by term.

\subsection{The combinator structural laws}

The data-algebra combinators of \cref{ch:combinators} carry structural laws that
organize their whole family. These are restated here as the reference catalogue; the
chapter derives them.

\begin{theorem}[\texttt{linear\_combination} is a monoid homomorphism]
\label{thm:lincomb}
\code{linear\_combination} carries the monoid of term-lists under concatenation to
the monoid of tensors under addition:
\begin{align}
  \textsf{linear\_combination}\ [\,] &\;=\; \mathbf{0}
    & \text{(empty list $\mapsto$ unit)} \label{eq:lcempty}\\
  \textsf{linear\_combination}\ (a \mathbin{+\!\!+} b)
    &\;=\; \textsf{linear\_combination}\ a + \textsf{linear\_combination}\ b
    & \text{(concatenation-homomorphism)} \label{eq:lcconcat}\\
  \textsf{linear\_combination}\ ((\kappa a, t):\mathit{rest})
    &\;=\; \textsf{linear\_combination}\ ((a, \kappa t):\mathit{rest})
    & \text{(scalar absorption)} \label{eq:lcabsorb}\\
  \textsf{linear\_combination}\ ((0, t):\mathit{rest})
    &\;=\; \textsf{linear\_combination}\ \mathit{rest}
    & \text{(zero-term drop)} \label{eq:lcdrop}
\end{align}
\end{theorem}

\noindent The concatenation-homomorphism \eqref{eq:lcconcat} is the structural fact
of \cref{ch:combinators}: it is what makes the four fixed-arity BLAS routines
(\code{scal}, \code{axpy}, \code{axpby}, \code{axpbypcz}) \emph{one} combinator,
since a longer term-list is the concatenation of shorter ones. The empty list maps to
the zero tensor \eqref{eq:lcempty}; a scalar slides freely between coefficient and
tensor \eqref{eq:lcabsorb}; a zero-coefficient term drops out \eqref{eq:lcdrop}. Over
an \emph{exact} field the sum is also permutation-invariant---but that last one is a
floating-point caveat, deferred to \cref{sec:nonlaws} as the load-bearing non-law it
becomes under IEEE-754.

\begin{theorem}[\texttt{inner\_product} structural laws]
\label{thm:ip}
\code{inner\_product} is the scalar-producing dual: a monoid homomorphism from
shape-concatenated tensors to scalar addition, bilinear in its arguments.
\begin{align}
  \textsf{inner\_product}\ (\vx_1 \mathbin{+\!\!+} \vx_2)\ (\vy_1 \mathbin{+\!\!+} \vy_2)
    &\;=\; \textsf{inner\_product}\ \vx_1\ \vy_1 + \textsf{inner\_product}\ \vx_2\ \vy_2
    \label{eq:ipsplit}\\
  \textsf{inner\_product}\ \mathbf{0}\ \vy
    &\;=\; \textsf{inner\_product}\ \vx\ \mathbf{0} \;=\; 0
    \label{eq:ipzero}
\end{align}
\end{theorem}

\noindent The split-additivity \eqref{eq:ipsplit} is the law that licenses blocked
and parallel evaluation of an inner product: split the operands, reduce the pieces
independently, add the partial sums. It is the formal twin of
\code{linear\_combination}'s concatenation law, only collapsing to a scalar rather
than a tensor. \Cref{fig:catalogue} gathers all of \cref{sec:coplaws} into a single
reference table.

\begin{figure}[t]
\centering
\renewcommand{\arraystretch}{1.35}
\setlength{\tabcolsep}{6pt}
\begin{tabular}{@{}p{0.26\textwidth} p{0.40\textwidth} p{0.24\textwidth}@{}}
\toprule
\textbf{Operator} & \textbf{Law} & \textbf{Guarantees} \\
\midrule
\code{axpy} &
  $\textsf{axpy}\ \alpha\ \vv\ \mathbf{0} = \alpha\vv$;\ \
  $\textsf{axpy}\ 0\ \vv\ \vy = \vy$ &
  peephole drop of a dead scale/add \\
\code{dot} &
  $\textsf{dot}\ \vv\ \vw = \textsf{dot}\ \vw\ \vv$ (symmetry) &
  operand order free (math, not bits) \\
\code{dot} &
  $\textsf{dot}\,(\textsf{axpy}\,\alpha\,\vv\,\vw)\,\vx
   = \alpha\,\textsf{dot}\,\vv\,\vx + \textsf{dot}\,\vw\,\vx$ &
  distribute over scaled sums \\
\code{matvec} &
  $\textsf{matvec}\,A\,(\alpha\vv) = \alpha\,\textsf{matvec}\,A\,\vv$ &
  pull scalars through apply \\
\code{matvec} &
  $\textsf{matvec}\,(A{+}B)\,\vv = \textsf{matvec}\,A\,\vv + \textsf{matvec}\,B\,\vv$ &
  assemble operator as a sum \\
\code{matvec} &
  $(\textsf{matvec}\,A)\!\circ\!(\textsf{matvec}\,B) = \textsf{matvec}\,(A\!\circ\! B)$ &
  fuse two applies into one \\
\code{linear\_combination} &
  $\textsf{lc}\,(a {+}\mkern-3mu{+}\, b) = \textsf{lc}\,a + \textsf{lc}\,b$ &
  split / chunk / reorder the sum (exact) \\
\code{inner\_product} &
  $\textsf{ip}\,(\vx_1{+}\mkern-3mu{+}\vx_2)\,(\vy_1{+}\mkern-3mu{+}\vy_2)
   = \textsf{ip}\,\vx_1\,\vy_1 + \textsf{ip}\,\vx_2\,\vy_2$ &
  blocked / parallel reduction (exact) \\
\bottomrule
\end{tabular}
\caption[The operator-law catalogue]{The algebraic operator laws of
\cref{sec:coplaws} as a reference table. Each row gives an operator, an equational
law, and the rewrite it permits. The two combinator homomorphism laws
(\code{linear\_combination}, \code{inner\_product}) are marked ``exact''---they hold
over an idealized field but become floating-point caveats once the reduction order
is fixed, the subject of \cref{sec:nonlaws} and \cref{app:numerics}. The
\code{matvec} composition law is the one with the most downstream leverage: it is the
algebraic license for assembling and fusing operators in the driven analysis.}
\label{fig:catalogue}
\end{figure}

\section{The laws that do \emph{not} hold}
\label{sec:nonlaws}

We arrive at the appendix's most valuable content. The laws above tell you which
rewrites are safe. This section tells you which \emph{plausible} rewrites are
\emph{unsafe}---the ones that resemble a law closely enough to tempt a refactor, but
fail, and fail silently. Each is stated as the false equation it is, with the reason
it fails and the concrete damage assuming it would do. They are grouped by the
operator they tempt you to mishandle.

\subsection{Non-laws of \texttt{iterate\_while}}
\label{sec:nlloop}

The loop combinator of \cref{ch:fold} carries four traps, all variations on ``treat
the loop as if it were an ordinary algebraic object.''

\begin{description}[style=nextline, leftmargin=0pt, font=\normalfont\itshape]
  \item[Predicate hoisting is invalid.] One might hope to hoist the loop predicate
  out of the iteration---test it once, then run the body unconditionally:
  \begin{lstlisting}[style=l4style]
iterate_while a p f  =/=  if p a then (repeat f forever) else stop   -- WRONG
  \end{lstlisting}
  This is false. The predicate is re-evaluated on \emph{every} iteration, not once at
  entry. Hoisting it converts a bounded loop into an unbounded one driven by the
  \emph{initial} predicate value---if \code{p a} is true at entry, the hoisted form
  runs forever, because nothing re-tests it. \textbf{What breaks:} a solver that was
  converging in fifty steps now never terminates; the bug is a hang, not a wrong
  number, and it surfaces only on inputs where the predicate was initially true.
  \item[Loop fusion across iterations is invalid.] Two \code{iterate\_while}s in
  sequence are \emph{not} one \code{iterate\_while} with a merged predicate:
  \begin{lstlisting}[style=l4style]
(iterate_while a p1 f) then (iterate_while _ p2 f)  =/=  iterate_while a (p1 || p2) f
  \end{lstlisting}
  except in the special case where the two predicates agree at the handoff state.
  The two-phase fold's trajectory is the concatenation \code{traj1 ++ traj2} with the
  second phase \emph{restarting} from the first phase's final state; the flattened
  single fold is one continuous trajectory under one predicate. \textbf{What breaks:}
  this is precisely why GMRES is written as an \emph{outer} restart loop wrapped
  around an \emph{inner} \code{iterate\_while} (\cref{ch:gmres}), not as one flat
  loop---the restart \emph{is} the phase boundary, and fusing it away discards the
  restart's basis reset. Merging the two loops silently changes the algorithm from
  restarted GMRES to a single un-restarted Krylov method with a different convergence
  and memory profile.
  \item[\code{while} and \code{do-while} are not interchangeable.] The book's
  \code{iterate\_while} tests the predicate \emph{before} the body (a \code{while});
  the variant that tests \emph{after} (a \code{do-while}, which always runs the body
  at least once) is a \emph{different combinator}:
  \begin{lstlisting}[style=l4style]
iterate_while a p f  =/=  iterate_while_post a p f   -- differ when p a is false at entry
  \end{lstlisting}
  They disagree exactly when the predicate is false at entry: the \code{while} runs
  zero steps and returns \code{a}; the \code{do-while} runs one step anyway.
  \textbf{What breaks:} swapping one for the other off-by-ones the iteration count and,
  for a loop entered already-converged, performs one spurious step on a state that
  should have been returned untouched.
  \item[No identity element.] There is no carry \code{a\_id} for which
  \code{iterate\_while a\_id p f} is a unit for all \code{p}, \code{f}. The loop is a
  \emph{fold}, not a monoid operation; folds do not have a unit in the carry.
  \textbf{What breaks:} any refactor that tries to ``initialize the loop with a
  neutral state'' and expects it to be inert is resting on a structure that is not
  there.
\end{description}

\subsection{Non-laws of \texttt{krylov\_step}}
\label{sec:nlkrylov}

The Krylov step is the richest source of non-laws, because it \emph{looks} like a
composition of well-behaved linear primitives and is not---its scalar updates involve
divisions and its convergence flag a comparison, and those destroy the linearity and
composability the primitives individually enjoy. Each non-law below is a property the
\emph{primitives} have that the \emph{step} does not.

\begin{description}[style=nextline, leftmargin=0pt, font=\normalfont\itshape]
  \item[The primitive sequence does not commute.] The five primitive groups of a
  step---operator apply, auxiliary, iterate-update, scalar-update, output-readout---
  cannot be reordered. The dataflow is rigid: \code{apply} produces $\vw$, the
  \code{axpy} reads it, the \code{dot} reads it. \textbf{What breaks:} reorder the
  apply past the \code{dot} that consumes its output and you read uninitialized or
  stale data; the result is not a different valid ordering, it is wrong.
  \cref{fig:noncommute} traces two orderings of a reduction to bit-different results.
  \item[No fold-merge / associativity.] Folding a Krylov step over predicate $p_1$
  and then over $p_2$ is \emph{not} folding over a merged predicate
  \lstinline[style=l4style]|p1 || p2|. The inner fold's per-step outputs are invisible
  to the outer fold, and convergence predicates that depend on monotonic-loss
  properties do not compose. \textbf{What breaks:} the same restart-flattening hazard
  as the loop's, now at the step level---it is why restart logic is an outer loop
  (\cref{ch:gmres}), not a flattened single fold.
  \item[No identity element.] There is no Krylov state \code{K\_id} with
  \lstinline[style=l4style]|krylov_step op K_id = pure { krylov: K_id, ... }|. The
  tempting candidate---a zero step length $\alpha = 0$---is not identity but
  \emph{breakdown}: in CG, $\alpha = 0$ is a division-by-zero-adjacent degeneracy, not
  a no-op. \textbf{What breaks:} treating a stalled step as an identity hides a
  breakdown the solver must detect and report.
  \item[No linearity in any argument.] The step is \emph{not} linear in its state:
  \[
    \textsf{krylov\_step}\ op\ (\alpha K_1 + \beta K_2)
      \;\neq\; \alpha\,\textsf{krylov\_step}\ op\ K_1
             + \beta\,\textsf{krylov\_step}\ op\ K_2.
  \]
  It is built from linear primitives, but their composition with \code{dot} and the
  scalar \emph{divisions} of the coefficient updates ($\alpha = r\!\cdot\! r / p\!\cdot\! Ap$)
  is non-linear, and the convergence comparison is not even continuous.
  \textbf{What breaks:} any attempt to ``superpose'' two solves into one by linearity
  computes nonsense; the step is an affine-ish iteration, not a linear map.
  \item[Bit-non-determinism across orthogonalization variants.] Switching the
  orthogonalization scheme---modified Gram--Schmidt (MGS), classical (CGS), or
  twice-iterated (CGS2)---produces \emph{mathematically equivalent} but
  \emph{bit-distinct} trajectories. \textbf{What breaks:} a regression test that pins
  exact bits will fail across a variant switch even though the math is unchanged; this
  is a load-bearing numerical choice (\cref{ch:bigidea}'s transparent-versus-load-bearing
  distinction), not a free refactor.
  \item[Form-A / Form-B equivalence is \emph{not} a monad-law rewrite.] The
  first-iteration-unrolled Form B of a Krylov step produces
  iteration-for-iteration-identical trajectories to the rolled Form A---but the two
  are \emph{not} related by the monad laws of \cref{sec:cmonad}. The rotation between
  them is a \emph{structural} rewrite (it drops a state field and threads a closure
  argument), not a $\beta$-reduction or a \code{>>=}-associativity step.
  \textbf{What breaks:} citing ``the monad laws'' to justify the unrolling is a
  category error; the equivalence holds, but it must be argued structurally, and a
  tool that only knows the monad laws cannot verify it.
\end{description}

\begin{figure}[t]
\centering
\begin{tikzpicture}[font=\footnotesize]
  % ===== ordering A =====
  \begin{scope}
    \node[font=\small\bfseries, text=simBlue] at (0,2.45) {ordering A};
    \node[data, minimum width=11mm] (a0) at (-1.4,1.5) {$c_1$};
    \node[opbox, minimum width=9mm] (aA) at (0.2,1.5) {$+$};
    \node[data, minimum width=11mm] (a1) at (1.8,1.5) {$c_2$};
    \draw[flow] (a0)--(aA); \draw[flow] (a1)--(aA);
    \node[opbox, minimum width=9mm] (aB) at (0.2,0.4) {$+$};
    \node[data, minimum width=11mm] (a2) at (1.8,0.4) {$c_3$};
    \draw[flow] (aA)--(aB); \draw[flow] (a2)--(aB);
    \node[product, minimum width=22mm] (ar) at (0.2,-0.7)
      {$(c_1\!\oplus\! c_2)\!\oplus\! c_3$};
    \draw[flow] (aB)--(ar);
    \node[font=\scriptsize\ttfamily, text=simRust] at (0.2,-1.5) {= 1.0000003};
  \end{scope}
  \draw[simGray!50, thick] (3.3,-1.8) -- (3.3,2.6);
  % ===== ordering B =====
  \begin{scope}[xshift=6.6cm]
    \node[font=\small\bfseries, text=simGreen] at (0,2.45) {ordering B};
    \node[data, minimum width=11mm] (b0) at (-1.4,1.5) {$c_2$};
    \node[opbox, minimum width=9mm] (bA) at (0.2,1.5) {$+$};
    \node[data, minimum width=11mm] (b1) at (1.8,1.5) {$c_3$};
    \draw[flow] (b0)--(bA); \draw[flow] (b1)--(bA);
    \node[opbox, minimum width=9mm] (bB) at (0.2,0.4) {$+$};
    \node[data, minimum width=11mm] (b2) at (1.8,0.4) {$c_1$};
    \draw[flow] (bA)--(bB); \draw[flow] (b2)--(bB);
    \node[product, minimum width=22mm] (br) at (0.2,-0.7)
      {$c_1\!\oplus\!(c_2\!\oplus\! c_3)$};
    \draw[flow] (bB)--(br);
    \node[font=\scriptsize\ttfamily, text=simRust] at (0.2,-1.5) {= 1.0000002};
  \end{scope}
  \node[font=\scriptsize\itshape, text=simGray, align=center] at (3.3,-2.3)
    {same operands, same exact sum --- different rounding, different last bit};
\end{tikzpicture}
\caption[Non-commutativity of a reduction: two orderings, different last bits]{Two
reduction orderings of the same three floating-point operands $c_1, c_2, c_3$, with
$\oplus$ the rounding floating-point add. Ordering A accumulates
$(c_1 \oplus c_2) \oplus c_3$; ordering B accumulates $c_1 \oplus (c_2 \oplus c_3)$.
The \emph{exact} sum is identical and order-independent---real addition is
associative---but each rounding step discards low bits differently, so the two orders
return results that differ in their last bit (the schematic values
\code{1.0000003} versus \code{1.0000002}). This is why \code{dot}, \code{nrm2}, and
\code{linear\_combination} are \emph{not} bit-associative: the reduction tree is part
of the algorithm, not a free choice (\cref{sec:cfp}, \cref{app:numerics}).}
\label{fig:noncommute}
\end{figure}

\subsection{Floating-point non-laws of the reductions}
\label{sec:cfp}

The most consequential non-laws in the entire book are the ones that hold in
mathematics and fail in floating point. The reduction verbs---\code{dot},
\code{nrm2}, \code{inner\_product}, and the summation inside
\code{linear\_combination}---are over an \emph{exact} field associative and
commutative, so the catalogue of \cref{sec:coplaws} stated their split and
permutation laws as exact. Under IEEE-754 those laws become \emph{caveats}.

\begin{description}[style=nextline, leftmargin=0pt, font=\normalfont\itshape]
  \item[Reductions are not bit-associative.] Floating-point $\oplus$ is not
  associative: $(a \oplus b) \oplus c \neq a \oplus (b \oplus c)$ in general, because
  each step rounds. So the value of \code{dot v w} or
  \lstinline[style=l4style]|linear_combination pairs| \emph{depends on the order in
  which the partial sums are accumulated}---the reduction tree. The mathematical value
  is order-agnostic; the computed value is not. \cref{fig:noncommute} is this fact in
  three operands.
  \item[Permutation-invariance of \code{linear\_combination} \eqref{eq:lcconcat} is
  exact-only.] Reordering the term list leaves the \emph{exact} sum unchanged but
  perturbs the rounding. The framework therefore pins the left-to-right
  \code{foldl} order as \emph{canonical} and treats any reordering as a load-bearing
  numerical change, not a free rewrite.
  \item[Split-additivity of \code{inner\_product} \eqref{eq:ipsplit} licenses blocked
  evaluation only up to rounding.] Splitting the operands and adding the partial
  inner products gives the \emph{exact} answer but a \emph{bit-different} one from the
  unsplit reduction, because the partial-sum tree differs.
  \item[Cauchy--Schwarz can fail by ULPs.] $\abs{\inner{\vx}{\vy}} \le \norm{\vx}\norm{\vy}$
  holds mathematically but can be violated by unit-in-last-place amounts in floating
  point, because the two sides are reduced by different trees.
\end{description}

\textbf{What breaks if you assume these are exact laws:} bit-reproducibility. A
solver refactored to ``reorder the reduction for cache locality'' or ``split the dot
product across two threads and add the halves'' computes a \emph{mathematically
equal} but \emph{bit-different} residual. For a single step that is invisible; over
thousands of iterations the bit-difference compounds through the convergence
trajectory, and two builds that should be identical now disagree on iteration counts
and final digits. This is the precise mechanism by which \emph{a refactor that looks
valid silently changes the last bits}, and it is the reason the book classifies
reduction order as a load-bearing numerical trick (\cref{ch:bigidea}) rather than a
transparent one. The full treatment---error bounds, compensated summation, the cost
of pinning a tree---is \cref{app:numerics}.

\begin{keyidea}
The floating-point non-laws are the sharpest edge in this appendix. An exact-field
law (associativity of $+$, permutation-invariance of a sum, Cauchy--Schwarz) is
\emph{true in mathematics and false in the machine}. Assuming the exact law when
refactoring a reduction does not produce a wrong answer you can see---it produces a
\emph{bit-different} answer that is mathematically correct and reproducibly wrong, and
the divergence only becomes visible after it has compounded through a long iteration.
Reduction order is part of the algorithm. Treat \code{dot}, \code{nrm2}, and the
\code{linear\_combination} sum as having a \emph{pinned tree}, not a free one, whenever
bit-reproducibility matters (\cref{app:numerics}).
\end{keyidea}

\section{Law versus non-law: the reference card}
\label{sec:ccard}

\Cref{fig:lawcard} is the two-column summary the rest of this appendix exists to
support: on the left, the rewrites you \emph{can} apply; on the right, the
look-alikes you \emph{cannot}. Consult it before any refactor of solver code.

\begin{figure}[t]
\centering
\renewcommand{\arraystretch}{1.3}
\setlength{\tabcolsep}{6pt}
\begin{tabular}{@{}p{0.46\textwidth} p{0.46\textwidth}@{}}
\toprule
\textbf{You CAN rewrite\ \ (laws)} & \textbf{You CANNOT rewrite\ \ (non-laws)} \\
\midrule
flatten / split / inline \code{do}-blocks \eqref{eq:assoc} &
  hoist a loop predicate out of \code{iterate\_while} \\
drop a \code{return} immediately consumed \eqref{eq:leftid} &
  fuse two sequential loops into one merged-predicate loop \\
fold get-apply-put into one \code{modify} \eqref{eq:modifydef} &
  swap \code{while} for \code{do-while} \\
fuse two adjacent \code{modify}s to $f\!\circ\! g$ \eqref{eq:modifyfuse} &
  collapse two \code{modify}s to the last one (that is a \code{put} law) \\
delete an inert get-then-put \eqref{eq:getput} &
  reorder the primitive sequence inside a \code{krylov\_step} \\
distribute \code{dot} over a scaled sum \eqref{eq:dotbilin} &
  treat $\alpha=0$ as a step identity (it is breakdown) \\
fuse two \code{matvec}s into one \eqref{eq:mvcomp} &
  superpose two Krylov solves by ``linearity'' \\
split a \code{linear\_combination} / \code{inner\_product} \emph{exactly} &
  reorder a reduction and expect \emph{bit}-identity \\
\bottomrule
\end{tabular}
\caption[Law versus non-law reference card]{The two-column reference card. The
left column lists the safe rewrites of \cref{sec:cmonad}--\cref{sec:coplaws}, each
keyed to its law; the right column lists the look-alike rewrites of
\cref{sec:nonlaws} that are \emph{not} laws. The pairing is deliberate: each
right-column entry resembles a left-column one closely enough to tempt the
substitution. The distinction between fusing two \code{modify}s (a law,
left) and collapsing two writes to the last (a \code{put} law misapplied to
\code{modify}, right) is the kind of near-miss this card exists to catch. When a
refactor is not on the left, assume it is on the right until proven otherwise.}
\label{fig:lawcard}
\end{figure}

\summarysection
This appendix is the reference catalogue of the rewrites the calculus licenses, in
two halves. The \emph{monad laws}---left identity \eqref{eq:leftid}, right identity
\eqref{eq:rightid}, associativity \eqref{eq:assoc}---govern \code{>>=} and
\code{return} and make \code{do}-block refactoring sound: you may erase a consumed
\code{return}, drop a trailing one, and flatten, split, or inline sub-blocks freely
(\cref{sec:cmonad}). The \emph{state-effect laws} \eqref{eq:modifydef}--\eqref{eq:putget}
are the algebra of \code{get}/\code{put}/\code{modify}: \code{modify} is get-apply-put,
adjacent \code{modify}s fuse to a composition, a get-then-put is inert, a later
\code{put} wins, and a get after a put sees the put---with the load-bearing asymmetry
that \code{put}s overwrite (so an earlier one dies) while \code{modify}s read first
(so they compose, not collapse) (\cref{sec:cstate}). The \emph{operator laws}
(\cref{sec:coplaws}) are the equational reference for the tensor primitives:
\code{axpy} corner cases, \code{dot} symmetry and bilinearity, \code{matvec}
linearity and composition, and the combinator homomorphisms that make
\code{linear\_combination} and \code{inner\_product} their families' entries. The
second half---the \emph{non-laws} (\cref{sec:nonlaws})---is the load-bearing content:
predicate hoisting and loop fusion are invalid; \code{while} and \code{do-while}
differ; a Krylov step does not commute, merge, have an identity, or stay linear, and
its Form-A/Form-B equivalence is structural, not a monad-law rewrite; and---sharpest
of all---the reductions are \emph{not} bit-associative, so reordering a \code{dot} or
a \code{linear\_combination} sum is a load-bearing numerical change, not a free one
(\cref{sec:cfp}). The reference card of \cref{sec:ccard} pairs each safe rewrite with
its dangerous look-alike. The governing lesson is the one the \cref{sec:cread} key
idea states and \cref{sec:cfp} sharpens: \textbf{the non-laws are as important as the
laws, because they are why a refactor that looks valid can silently change the last
bits.}

\furtherreading
The monad laws and the \code{get}/\code{put}/\code{modify} algebra are standard
functional-programming material; Wadler's tutorial derives them from first principles
and is the readable canonical source for why they make stateful-looking pure code
sound to refactor~\citep{wadler1995monads}, while the precise desugaring of
\code{do}-notation into \code{>>=}---on which the associativity law's
\code{do}-block reading rests---is documented in the language report Peyton Jones
edits~\citep{peytonjones2003}. The floating-point non-laws of \cref{sec:cfp}---the
non-associativity of summation and its consequences for iterative solvers---are
treated at length in Saad's standard reference on iterative methods, which is also
the home of the Krylov convergence theory the step's non-linearity and
non-composability (\cref{sec:nlkrylov}) ultimately rest on~\citep{saad2003}. The
reduction-order error analysis those non-laws demand is the subject of the companion
\cref{app:numerics}; the formal reduction rules these equational laws sit atop are in
\cref{app:calculus}.
